% GENERATED FILE. Edit canonical lessons, then run npm run generate:books.
% canonical-source-hash: fnv1a64:f0e5e68c28fb336b
% canonical-lessons: FA-C352-bastan, FA-C352-bastani, FA-C352-tamaddon, FA-C352-emperaturi, FA-C352-selsele

\chapter{Ancient times}
\label{ch:fa-ancient-times}
\hlchaptermodality{352}

\begin{chapteropening}
\textbf{By the end of this chapter:} \emph{I can talk about antiquity, ancient things, a civilisation, an empire and a dynasty.}

Complete the last lesson of chapter 352: I can talk about antiquity, ancient things, a civilisation, an empire and a dynasty.
\end{chapteropening}

\section[\texorpdfstring{\fa{باستان} (bâstân)}{باستان (bâstân)}]{\fa{باستان} (bâstân) — antiquity, ancient times}

\label{lesson:FA-C352-bastan}



\begin{quote}
Before the new one: say the Persian for to defeat, then the Persian for to conquer.
\end{quote}

\subsection*{You'll want to know: \fa{باستان}}
\textbf{\fa{باستان}} (\emph{bâstân}) — \textquotedblleft{}antiquity, ancient times\textquotedblright{}.

\begin{quote}
\fa{این} \fa{بنا} \fa{از} \fa{زمان} \fa{باستان} \fa{مانده} \fa{است}.
\end{quote}

\begin{quote}
\emph{In banâ az zamân-e bâstân mânde ast.}
\end{quote}

\begin{quote}
This building has survived from ancient times.
\end{quote}

\emph{Bâstânshenâs} \textquotedblleft{}an archaeologist\textquotedblright{}.

\begin{cousinweb}
\emph{bâstân} is a native Persian word that EtymDB glosses \textquotedblleft{}ancient era\textquotedblright{} without a parent.
\end{cousinweb}

\begin{grammarlens}[title={What you've built}]
One more word for telling, explaining and describing at B1.
\end{grammarlens}

\subsection*{Your turn}
\begin{itemize}
\raggedright
  \item \emph{Say it:} \emph{bâstân}
  \item \emph{Say it:} \emph{bâstân}, once more
  \item \emph{Say it:} a sentence of your own with \emph{bâstân}
  \item \emph{Recall:} say the Persian for to defeat, then the Persian for to conquer, then say \emph{bâstân} again
\end{itemize}

\begin{tcolorbox}[breakable,colback=teal!4,colframe=teal!35!black,title={Before you move on}]
What does \fa{باستان} mean? (\textquotedblleft{}Antiquity, ancient times\textquotedblright{}.) Say it once more.
\end{tcolorbox}
\section[\texorpdfstring{\fa{باستانی} (bâstâni)}{باستانی (bâstâni)}]{\fa{باستانی} (bâstâni) — ancient}

\label{lesson:FA-C352-bastani}



\begin{quote}
Before the new one: say the Persian for to conquer, then the Persian for antiquity.
\end{quote}

\subsection*{You'll want to know: \fa{باستانی}}
\textbf{\fa{باستانی}} (\emph{bâstâni}) — \textquotedblleft{}ancient\textquotedblright{}.

\begin{quote}
\fa{در} \fa{موزه} \fa{سکه} \fa{های} \fa{باستانی} \fa{دیدیم}.
\end{quote}

\begin{quote}
\emph{Dar muze sekke-hâ-ye bâstâni didim.}
\end{quote}

\begin{quote}
We saw ancient coins in the museum.
\end{quote}

\emph{Varzesh-e bâstâni} is the traditional strength training of the \emph{zurkhâne}.

\begin{cousinweb}
\emph{bâstâni} is built from Persian \emph{bâstân} \textquotedblleft{}antiquity\textquotedblright{}, a native word, with the adjective suffix \emph{-i}.
\end{cousinweb}

\begin{grammarlens}[title={What you've built}]
One more word for telling, explaining and describing at B1.
\end{grammarlens}

\subsection*{Your turn}
\begin{itemize}
\raggedright
  \item \emph{Say it:} \emph{bâstâni}
  \item \emph{Say it:} \emph{bâstâni}, once more
  \item \emph{Say it:} a sentence of your own with \emph{bâstâni}
  \item \emph{Recall:} say the Persian for to conquer, then the Persian for antiquity, then say \emph{bâstâni} again
\end{itemize}

\begin{tcolorbox}[breakable,colback=teal!4,colframe=teal!35!black,title={Before you move on}]
What does \fa{باستانی} mean? (\textquotedblleft{}Ancient\textquotedblright{}.) Say it once more.
\end{tcolorbox}
\section[\texorpdfstring{\fa{تمدن} (tamaddon)}{تمدن (tamaddon)}]{\fa{تمدن} (tamaddon) — a civilisation}

\label{lesson:FA-C352-tamaddon}



\begin{quote}
Before the new one: say the Persian for antiquity, then the Persian for ancient.
\end{quote}

\subsection*{You'll want to know: \fa{تمدن}}
\textbf{\fa{تمدن}} (\emph{tamaddon}) — \textquotedblleft{}a civilisation\textquotedblright{}.

\begin{quote}
\fa{تمدن} \fa{در} \fa{کنار} \fa{رودخانه} \fa{ها} \fa{شروع} \fa{شد}.
\end{quote}

\begin{quote}
\emph{Tamaddon dar kenâr-e rudkhâne-hâ shoru' shod.}
\end{quote}

\begin{quote}
Civilisation began on the banks of rivers.
\end{quote}

Plural \emph{tamaddon-hâ}. The Arabic word is linked with \emph{madina} \textquotedblleft{}city\textquotedblright{}: living in cities.

\begin{cousinweb}
\emph{tamaddon} is borrowed from Arabic \emph{tamaddun}, \textquotedblleft{}civilisation\textquotedblright{} (EtymDB), a word usually linked with \emph{madīna} \textquotedblleft{}city\textquotedblright{}.
\end{cousinweb}

\begin{grammarlens}[title={What you've built}]
One more word for telling, explaining and describing at B1.
\end{grammarlens}

\subsection*{Your turn}
\begin{itemize}
\raggedright
  \item \emph{Say it:} \emph{tamaddon}
  \item \emph{Say it:} \emph{tamaddon}, once more
  \item \emph{Say it:} a sentence of your own with \emph{tamaddon}
  \item \emph{Recall:} say the Persian for antiquity, then the Persian for ancient, then say \emph{tamaddon} again
\end{itemize}

\begin{tcolorbox}[breakable,colback=teal!4,colframe=teal!35!black,title={Before you move on}]
What does \fa{تمدن} mean? (\textquotedblleft{}A civilisation\textquotedblright{}.) Say it once more.
\end{tcolorbox}
\section[\texorpdfstring{\fa{امپراتوری} (emperâturi)}{امپراتوری (emperâturi)}]{\fa{امپراتوری} (emperâturi) — an empire}

\label{lesson:FA-C352-emperaturi}



\begin{quote}
Before the new one: say the Persian for ancient, then the Persian for a civilisation.
\end{quote}

\subsection*{You'll want to know: \fa{امپراتوری}}
\textbf{\fa{امپراتوری}} (\emph{emperâturi}) — \textquotedblleft{}an empire\textquotedblright{}.

\begin{quote}
\fa{این} \fa{امپراتوری} \fa{صدها} \fa{سال} \fa{دوام} \fa{آورد}.
\end{quote}

\begin{quote}
\emph{In emperâturi sad-hâ sâl davâm âvard.}
\end{quote}

\begin{quote}
This empire lasted for hundreds of years.
\end{quote}

\emph{Emperâtur} \textquotedblleft{}an emperor\textquotedblright{} + \emph{-i}.

\begin{cousinweb}
\emph{emperâturi} is built on \emph{emperâtur} \textquotedblleft{}emperor\textquotedblright{}, which EtymDB gives as a borrowing of Latin \emph{imperator}, \textquotedblleft{}a commander-in-chief\textquotedblright{}, later the title of the Roman emperors (Lewis \& Short).
\end{cousinweb}

\begin{grammarlens}[title={What you've built}]
One more word for telling, explaining and describing at B1.
\end{grammarlens}

\subsection*{Your turn}
\begin{itemize}
\raggedright
  \item \emph{Say it:} \emph{emperâturi}
  \item \emph{Say it:} \emph{emperâturi}, once more
  \item \emph{Say it:} a sentence of your own with \emph{emperâturi}
  \item \emph{Recall:} say the Persian for ancient, then the Persian for a civilisation, then say \emph{emperâturi} again
\end{itemize}

\begin{tcolorbox}[breakable,colback=teal!4,colframe=teal!35!black,title={Before you move on}]
What does \fa{امپراتوری} mean? (\textquotedblleft{}An empire\textquotedblright{}.) Say it once more.
\end{tcolorbox}
\section[\texorpdfstring{\fa{سلسله} (selsele)}{سلسله (selsele)}]{\fa{سلسله} (selsele) — a dynasty; a chain}

\label{lesson:FA-C352-selsele}



\begin{quote}
Before the new one: say the Persian for a civilisation, then the Persian for an empire.
\end{quote}

\subsection*{You'll want to know: \fa{سلسله}}
\textbf{\fa{سلسله}} (\emph{selsele}) — \textquotedblleft{}a dynasty; a chain\textquotedblright{}.

\begin{quote}
\fa{این} \fa{کاخ} \fa{را} \fa{پادشاهان} \fa{یک} \fa{سلسله} \fa{قدیمی} \fa{ساختند}.
\end{quote}

\begin{quote}
\emph{In kâkh râ pâdeshâhân-e yek selsele-ye qadimi sâkhtand.}
\end{quote}

\begin{quote}
The kings of an old dynasty built this palace.
\end{quote}

Its first meaning is \textquotedblleft{}a chain\textquotedblright{}: \emph{selsele-ye kuh-hâ} \textquotedblleft{}a mountain range\textquotedblright{}.

\begin{cousinweb}
\emph{selsele} is borrowed from Arabic \emph{silsila}, \textquotedblleft{}chain\textquotedblright{} (EtymDB): a chain of kings.
\end{cousinweb}

\begin{grammarlens}[title={What you've built}]
One more word for telling, explaining and describing at B1.
\end{grammarlens}

\subsection*{Your turn}
\begin{itemize}
\raggedright
  \item \emph{Say it:} \emph{selsele}
  \item \emph{Say it:} \emph{selsele}, once more
  \item \emph{Say it:} a sentence of your own with \emph{selsele}
  \item \emph{Recall:} say the Persian for a civilisation, then the Persian for an empire, then say \emph{selsele} again
\end{itemize}

\begin{tcolorbox}[breakable,colback=teal!4,colframe=teal!35!black,title={Before you move on}]
What does \fa{سلسله} mean? (\textquotedblleft{}A dynasty; a chain\textquotedblright{}.) Say it once more.
\end{tcolorbox}
